% Generated draft: requires mathematical and source review.
\documentclass[11pt]{article}
\usepackage[margin=25mm]{geometry}
\usepackage{fontspec}
\setmainfont{Latin Modern Roman}
\usepackage{amsmath,amssymb,mathtools,unicode-math}
\setmathfont{Latin Modern Math}
\usepackage{enumitem,needspace,xurl,hyperref}
\setlength{\parindent}{0pt}
\setlength{\parskip}{5pt}
\setlength{\emergencystretch}{3em}
\newcommand{\N}{\mathbb{N}}
\newcommand{\Z}{\mathbb{Z}}
\newcommand{\Q}{\mathbb{Q}}
\newcommand{\R}{\mathbb{R}}
\newcommand{\C}{\mathbb{C}}
\newcommand{\F}{\mathbb{F}}
\newcommand{\A}{\mathbb{A}}
\newcommand{\PP}{\mathbb{P}}
\newcommand{\E}{\mathbb{E}}
\newcommand{\eps}{\varepsilon}
\newcommand{\dd}{\,\mathrm{d}}
\begin{document}
\section*{UnsolvedMath ID: 11000211}
\textbf{Source problem code:} catalogue entry
\subsection{Definitions and mathematical statement}
A translation surface \((X,\omega)\) consists of a compact Riemann surface and a nonzero holomorphic \(1\)-form; integration of \(\omega\) gives translation charts away from its zeros. Its Veech group is the derivative group of orientation-preserving affine automorphisms, viewed in \(\operatorname{SL}(2,\R)\). The surface is Veech when this group is a lattice, meaning finite covolume. For a zero-order partition \(k_1+\cdots+k_r=2g-2\), the corresponding stratum consists of forms with precisely these zero multiplicities, with the genus-one empty-zero convention.

The problem is to characterize all Veech surfaces in every stratum of every genus. The source does not prescribe an equivalence relation, normalization or format of classification; in particular it does not require a finite list or restrict to arithmetic surfaces.

It reports the genus-one case as trivial, genus-two solutions by K. Calta [Ca] and C. McMullen [McM2], and classification in [KnSm] and Puchta [Pu] of acute rational billiard triangles whose unfolded translation surfaces are Veech. Rational means that all angles are rational multiples of \(\pi\). It also cites [SmWe] for a Veech-surface criterion in terms of areas of triangles embedded in \((X,\omega)\). The exact triangle criterion and the scope of the stated classifications are not printed; no formula, exhaustive higher-genus list or reduction to arithmeticity is invented.
\subsection{Short English statement}
Classify the translation surfaces with lattice Veech group in every zero stratum and genus, preserving the reported genus-two, billiard-triangle and embedded-triangle-area results as context.
\subsection{Sources}
\begin{itemize}
\item \emph{Problems on Mapping Class Groups and Related Topics}, edited by Farb, Problem 4, ``Characterization of Veech surfaces,'' supplied PDF page 259. \url{https://www.math.uchicago.edu/~farb/papers/mcgbook.pdf}. Recorded access 2026-07-29; extraction pdf-font-structure.
\item K. Calta [Ca], C. McMullen [McM2], [KnSm], Puchta [Pu], and [SmWe], with roles described in the source. Complete bibliography entries and precise theorem statements are absent.
\item Local UnsolvedMath registry, \texttt{data/unsolvedmath/problems.json}, permanent ID \(11000211\), exact source code \texttt{catalogue entry}; supplied OPEN-TRIAGE summary mentions partially open Veech classification without further exact citations.
\end{itemize}
\end{document}
